\lesson{1}{Nov 15 2021 Mon (08:01:30)}{Polynomial Long Division}{Unit 3}

\subsubsection*{Steps to Dividing Polynomials}

\begin{enumerate}
    \item Divide
    \item Multiply
    \item Subtract
    \item Bring down
    \item Repeat
\end{enumerate}

\subsubsection*{Missing terms}

When either polynomial does not have all terms according to the descending powers on the variable, it is necessary to add in terms with $0$ coefficients.

\begin{example}[Dividing Polynomials]
    \[ \polylongdiv{x^2 + 3x + 2}{x + 1} \]
    
    Let me quickly go over everything:
    
    \begin{enumerate}
        \item First part, I'm multiplying $x$ by $x$ to get $x^2$.
        \item Second part, I'm also multiplying $1$ by $2$ to get $2$ because every time we do something to one, we must do it to all.
        \item Third part, I'm moving the answer we get ($x^2 + 2$) down. Now, I also distributed a negative ($-$), because I'm subtracting.
        \item Then, I repeat. I move the $2$ down next to the result I got, which was $2x$.
        \item Then, I multiply $x$ by $2$, which gets me $2x$.
        \item Since I multiplied $x$ by $2$, I have to do the same on the $1$. In the end, it brought me to $2x + 1$. But, I wanted to subtract the entire thing from the other answer, so I negated the whole equation.
        \item After subtracting, you are left with a $0$, meaning there is no remainder.
    \end{enumerate}
\end{example}

\subsubsection*{Remainders}

If the quotient has no remainder (a remainder of $0$), then the divisor is said to be a factor of the dividend.

\newpage
